\documentclass[11pt,a4paper]{article}
\usepackage[T1]{fontenc}
\usepackage{lmodern,amsmath,amssymb,amsthm,mathtools,booktabs,microtype}
\usepackage[margin=27mm]{geometry}
\usepackage[colorlinks=true,linkcolor=blue,citecolor=blue,urlcolor=blue]{hyperref}
\setlength{\emergencystretch}{2em}
\newtheorem{theorem}{Theorem}[section]
\newtheorem{lemma}[theorem]{Lemma}
\newtheorem{proposition}[theorem]{Proposition}
\newcommand{\Q}{\mathbb Q}\newcommand{\Z}{\mathbb Z}
\DeclareMathOperator{\Tr}{Tr}
\DeclareMathOperator{\Nm}{N}
\title{Scalar precision two with an eighth-order elliptic comparison}
\author{}\date{17 November 2026}
\begin{document}\maketitle
\begin{abstract}
We give a complete revision of the general elliptic source in which
the final scalar tests have order two. An eighth-order logarithmic
truncation supplies sixth-order comparisons in the first formal
neighborhood. A squared Lucas parameter gives the growth needed for
soundness with the weaker scalar height bound. The root and auxiliary
point roles, two pole pairs, and one-coordinate finite-algebra encoder
are retained. Under the source's arithmetic premises the resulting
family defines the integers with twenty-five total variables. No
twenty-one-variable construction or numerical degree bound is asserted.
\end{abstract}

\section{The revised functions}

Retain the general normalized source of \cite{Reweighted}, including
its scalar \(\rho=\rho_L(X)\), integrality-reflecting label \(\beta(X)\),
consumed prime \(\ell\), exceptional set, two contact functions \(R_j\),
and unsquared scales \(a_0=1,a_1=1+\rho T^2\).
Here \(L\ge3\) is fixed and odd, and \(T\) is the rational trace of
the sum of the four root abscissas. Either scalar covered by that
source may be used. All group, contact-height, aggregate trace-height,
prime-pattern, and pole-realization premises keep their stated scope.

The curve and source point are
\[
 K=\Q(\sqrt2),\quad D=11-6\sqrt2,\quad E:y^2=x^3-Dx,
 \quad
 P=\left(\frac{57-28\sqrt2}{8},\frac{-660+417\sqrt2}{32}\right).
\]
The nine point roles are
\[
 A,S,T_1,Z,H,U_A,U_S,U_T,U_Z.
\]
Write the root labels as \((n,s,t,b)\), so the four root points are
\([n]P,[s]P,[t]P,[b]P\), and set
\[
 c=|s-1|,\qquad B_0=\max(|n|,|s|,|t|,|b|).
\]
Let \(Q_0=P\), \((Q_1,Q_2,Q_3,Q_4)=(A,S,T_1,Z)\),
\(U_0=H\), and \(V_i=U_i-Q_i\). Keep all native point and
difference guards.

For \(f=1/x\), \(z=-x/y\), put
\begin{equation}\label{eq:ell8}
 \ell_8(Q)=z(Q)\left(1-\frac25D f(Q)^2\right)
          =\frac{2D-5x(Q)^2}{5x(Q)y(Q)}.
\end{equation}
Define the four paired comparisons by
\[
 L_i=\frac{\ell_8(U_i)}{\ell_8(H)},\qquad
 L_i'=\overline{\frac{\ell_8(V_i)}{\ell_8(V_0)}}.
\]
The sixteen field tests are
\begin{equation}\label{eq:fieldtests}
 (f(U_i))_{i=0}^4,\quad
 (\overline{f(V_i)})_{i=0}^4,\quad
 (L_i-L_i')_{i=1}^4,\quad
 L_1-\beta,\quad \mathcal F(L_2^2,L_3,L_4),
\end{equation}
where
\[
 \mathcal F(a,t,b)=t^2-a(a-2)tb+(a-2)^2b^2-1.
\]
The ten reciprocal tests have order one. The final six tests have
orders \((1,6,6,6,1,6)\), respectively.

The six scalar tests remain the rational components of
\[
 f(H),\qquad f(S-P),\qquad z(S-P)/z(H),
\]
and have orders \(\kappa=(1,1,1,1,2,2)\).
In particular, the scalar ratio still uses \(z\), not \(\ell_8\).
For each contact block retain
\[
 u_j=\rho R_j,\quad w_j=\frac{u_j}{(1+u_j)^2},\quad
 G_{j,k}=R_j^{\kappa_k+1}\frac{R_jS_k}{1+R_jS_k},
\]
and use
\begin{equation}\label{eq:detector}
 W_j=\prod_{i=1}^{16}\left(1+
 \frac{C_i^2}{\delta_j^2w_j^{2a_i-1}}\right)
 \prod_{k=1}^6\left(1+\frac{G_{j,k}^2}{w_j^{2\kappa_k-1}}\right),
 \qquad b_j=\frac{4w_j}{\Tr(W_j^{-1})^2}.
\end{equation}
Both native pole pairs remain. All displayed divisions and syntactic
denominators are guarded; no unrestricted coordinate is introduced by
\eqref{eq:ell8}.

\section{Formal precision and scalar extraction}

\begin{lemma}\label{lem:formal}
At a good prime \(p>76176\), if a source point \(Q\) lies in the
first formal neighborhood and \(r=v(z(Q))\ge1\), then
\[
 \frac{\ell_8(Q)}{\log_E(Q)}\in1+p^{8r}\mathcal O.
\]
At every prepared real sign branch of a nonzero source multiple,
\(\ell_8(Q)\ne0\).
\end{lemma}
\begin{proof}
The exact source expansion is
\[
 \log_E(Q)=z(Q)\sum_{j\ge0}a_j(Df(Q)^2)^j,
 \quad a_0=1,\quad a_1=-2/5,\quad
 a_j=a_{j-1}\frac{4j-5}{4j+1}\quad(j\ge2).
\]
The omitted terms in \eqref{eq:ell8} start at \(j=2\), and
\(v(Df(Q)^2)=4r\). When \(4j+1<p\), the coefficient is
integral and its term has valuation at least \(8r\).
For all remaining terms, the denominator estimate used in the source
gives \(v(a_j)\ge-j\). Since then \(j\ge19044\),
\(4jr-j\ge3jr\ge8r\), and these valuations tend to infinity.
Thus the relative omitted tail has valuation at least \(8r\).
The full logarithm divided by \(z(Q)\) is a unit, proving the assertion.

At either real embedding a nonzero multiple of \(P\) has a positive
square abscissa, and its equation implies \(x(Q)^2>D\).
Consequently
\[
 1-\frac{2D}{5x(Q)^2}>\frac35>0.
\]
The old nonzero ordinate and abscissa guards therefore give the last claim.
The same argument applies to every prepared sign branch of every
guarded difference, since its index is again nonzero.
\end{proof}

\begin{lemma}\label{lem:local}
Let \(p\) be outside the retained exceptional set.
If \(\mu=v_p(\rho)\) is odd, every rational source solution satisfies
\[
 v_p(c)>\frac54|\mu|,\qquad v_p(c)\ge2.
\]
At a positive odd contact of \(R_j\), either all field tests have
their assigned orders at one place above \(p\), or \(p^2\mid c\).
\end{lemma}
\begin{proof}
The unchanged disjoint-contact argument gives a unit \(R_j\) at an
odd scalar valuation. Its physical weight has valuation \(|\mu|\).
The normalized trace lemma of \cite{Reweighted} applies to arbitrary
positive integer test orders. It gives the bounds \(|\mu|/4\) on
the first four scalar components and \(5|\mu|/4\) on the last two.
The unchanged formal-group isometry for the scalar ratio gives
\[
 v(z(S-P)/z(H))=v_p((s-1)/h),\qquad H=[h]P.
\]
Since \(h\) is a nonzero integer, the bound for \(c\) follows.
At a comparison contact with even scalar valuation, the same trace
lemma gives
\(v(C_i)>(a_i-3/4)|v_p(\rho R_j)|\ge a_i-3/4\).
Integral valuations give the assigned orders. An odd scalar valuation
instead gives \(p^2\mid c\).
\end{proof}

Suppose the comparison alternative holds. The reciprocal tests put
the five points in their first formal neighborhoods at the appropriate
conjugate places. Lemma~\ref{lem:formal} is therefore sufficient:
no second formal neighborhood is needed. The usual paired-ratio
argument first makes both exact index ratios integral, because at
least one of \(h,h-1\) is a unit. Relative precision eight then becomes
absolute precision eight. Hence
\[
 L_1\equiv n\pmod p,\qquad
 (L_2,L_3,L_4)\equiv(s,t,b)\pmod{p^6}.
\]
Writing \(\beta=u/v\) in lowest terms and
\(\Delta=\mathcal F(s^2,t,b)\), the resulting radical divisibilities are
\begin{equation}\label{eq:divisibilities}
 M_j^6\mid\Delta c^3,\qquad M_j\mid(u-nv)c.
\end{equation}
Indeed a scalar alternative gives \(p^6\mid c^3\), whereas a
comparison alternative gives \(p^6\mid\Delta\).

\section{Heights and soundness}

The scalar contact-height bound and Lemma~\ref{lem:local} give
\begin{equation}\label{eq:scalarheight}
 h(\beta)\le Cc^{7/6},\qquad
 h(\rho)\le Cc^{7/6}+C\log c+C.
\end{equation}
For the aligned scalar, apply the same estimate to its fixed affine
contact argument and absorb the fixed affine constants. The consumed
prime still has positive odd valuation at least \(L\); hence
\[
 v_\ell(c)>5L/4,\qquad c\longrightarrow\infty
 \quad\text{uniformly as }L\longrightarrow\infty.
\]
In particular \(2L\log\ell<(8/5)\log c\), which gives the logarithmic
term in \eqref{eq:scalarheight}.

The balanced contact-height estimate is unchanged:
\(h(t_j)\le C M_j^{7/6}\). The unsquared comparison scale still gives
\[
 2h(T)\le h(a_1)+h(\rho)+\log2,
 \qquad h(a_1)\le h(t_0)+h(t_1).
\]
Together with the aggregate trace-height theorem and
\eqref{eq:divisibilities}, these imply, when \(\Delta\ne0\),
\[
 B_0^2\le C(|\Delta|c^3)^{7/36}
               +Cc^{7/6}+C\log c+C.
\]
All constants are independent of \(L\) and the source tuple.
Since \(\Delta\) has degree six, \(|\Delta|\le5B_0^6\) for
\(B_0\ge2\), and \(c\le B_0+1\), the first exponent in \(B_0\)
is \((6+3)7/36=7/4<2\). The other polynomial exponent is
\(7/6<2\). The forced growth of \(c\) therefore gives
\(\Delta=0\) for every sufficiently large fixed odd \(L\).

Apply the integer Lucas growth lemma of \cite{Baseline} with parameter
\(a=s^2\). The point guards give \(b\ne0\), and sufficiently
large \(c\) gives \(|s|>3\). Thus
\[
 \log B_0\ge\log|t|>s^2-3
              \ge c(c-2)-2.
\]
If the input residual \(N=u-nv\) is nonzero, the second divisibility
in \eqref{eq:divisibilities} instead gives
\[
 B_0^2\le C(|N|c)^{7/6}+Cc^{7/6}+C\log c+C.
\]
The last three terms can be absorbed into \(B_0^2/2\) for large
\(c\). Using \(|N|\le2e^{h(\beta)}B_0\) and taking logarithms yields
\[
 \log B_0\le Cc^{7/6}+C\log c+C,
\]
which contradicts the preceding quadratic lower bound.
Therefore \(\beta=n\in\Z\), and the unchanged input label implies
\(X\in\Z\).

\section{Ordered completeness and every character guard}

Fix an integer input and any fixed odd \(L\ge3\).
Prepare the shared reference point, the scalar point \(S=[s]P\),
and the six scalar components in exactly the order of
\cite{General,Reweighted}. Choose \(s=c+1>1\) odd and congruent to
one modulo three. The scalar orders are lower, so the original
arbitrarily small scalar neighborhoods are sufficient.

For the Lucas labels use parameter \(a=s^2\). Let
\[
 U_0(a)=0,\quad U_1(a)=1,\quad
 U_{j+1}(a)=aU_j(a)-U_{j-1}(a),
\]
and choose
\[
 k=2(a-2),\qquad b=U_k(a)/(a-2),\qquad t=U_{k+1}(a).
\]
The congruence \(U_k(a)\equiv k\pmod{a-2}\) makes \(b\) integral,
and the norm identity gives \(\mathcal F(s^2,t,b)=0\).
As \(a\) is odd and \(a\equiv1\pmod3\), the mod-two sequence has
period \(0,1,1\), and \(k\equiv1\pmod3\). Thus \(t,b\) are positive
and odd. The input label remains \(n=\beta\). All roots and scales
are now fixed.

Temporarily put \(H=[h]P\), \(U_i=[hn_i]P\).
The new truncated ratios are continuous on the old prepared real
neighborhoods and fixed finite-place neighborhoods once their
denominators are guarded. Lemma~\ref{lem:formal} shows that the new
real denominator factors never vanish on any prepared sign branch.
The same finite avoidance and compact-neighborhood argument therefore
bounds every revised field test before prime selection.

Apply the unchanged simultaneous prime-pattern theorem, including the
original fixed avoidance (in particular the prime support of
\(1+\rho\)), denominator divisibility, positive large
slope, conjugate reduction-order conditions and extra divisibility
\(h\equiv0\pmod{\prod p_i}\). No deeper opposite-place congruence
is imposed. The compatible progression remains dense in the required
component of the compact cyclic closure, as in \cite{Unit}.

At a selected comparison prime, all reference and shifted points lie
in their first formal neighborhoods, so the reciprocal tests have
valuation at least two. The exact logarithmic ratios are the integers
\(n_i\). Lemma~\ref{lem:formal} consequently gives valuation at least
eight for every paired comparison and for the input target.
Because the exact Lucas residual vanishes, its polynomial comparison
also has valuation at least eight. For the order-six tests,
\[
 v\left(C_i^2/(\delta_j^2w_j^{11})\right)\ge16-11=5.
\]
For the order-one reciprocal tests the same lower bound is at least
\(4-1=3\). Hence all field factors are one modulo \(p^3\) at the
designated place.

The scalar normalization argument is unchanged. The original extra
divisibility makes the primary reference point have formal depth at
least two, while at the opposite place it reduces to \(P\).
The six rational scalar components therefore have nonnegative
valuation or valuation at most minus two. Thus \(R_jS_k\) never has
valuation zero at a fresh contact, and the normalizers are integral.
Now \(v(G_{j,k})\ge\kappa_k+1\), so every scalar second summand has
valuation at least three. The reciprocal trace is a unit, and the
fresh odd valuation of the pole input is positive.

At every fixed preparation place, \(R_j\to1\),
\(w_j\to\rho/(1+\rho)^2\), and \(\delta_j^{-1}\to0\) as the slope
tends to infinity. The new orders are fixed, the weights have a fixed
nonzero limit, and the revised tests are bounded on the already chosen
neighborhoods. Thus every old finite and real pole preparation remains
available. The same source factorization controls all odd weight primes.
Both pole inputs satisfy the retained converse hypotheses.

For the whole-character guards, root and difference signs still give
nonzero multiples \([\pm h\pm1]P\) and
\([\pm hn_i\pm n_i]P\), with \(h>2\) and \(n_i\ne0\).
Their new factors \(1-2Df^2/5\) are positive by
Lemma~\ref{lem:formal}. The scalar normalizers are precisely the old
ones, and their zero sets are avoided by the original analytic
avoidance argument after the compatible index refinement.
All revised detector factors are one plus a nonnegative square times
a positive coefficient at every real point-algebra character.
Their product and trace norms are therefore units.
The same density of the independent pole witnesses avoids every
remaining conjugate pole denominator and radicand.

\begin{theorem}
Under the retained arithmetic premises, the revised source defines
\(\Z\) over \(\Q\) for every sufficiently large fixed odd \(L\).
It is complete for every integer input and every fixed odd \(L\ge3\).
Its finite-algebra polynomial has twenty-five total variables:
one input, eighteen point coordinates, one slope, four pole coordinates,
and one encoder coordinate. No numerical determinant degree or smaller
defining-variable bound is asserted.
\end{theorem}
\begin{proof}
Soundness and ordered completeness have been proved above.
All added denominator factors lie in the existing point algebra and
are units at the prepared whole-character specializations. The guarded
finite-algebra encoder of \cite{Unit} therefore applies with unchanged
free coordinates and finite ranks.
\end{proof}

\begin{thebibliography}{9}\raggedright
\bibitem{General}\emph{A defining polynomial in twenty-five variables:
combined scalar and comparison conditions}, repository manuscript,
\nolinkurl{papers/few_variables25/few_variables25_merged_20261110_v19.tex}.
\bibitem{Reweighted}\emph{Quadratic reweightings of native elliptic detectors},
repository manuscript,
\nolinkurl{papers/pole_parameters/reweighted_elliptic_detectors_20261116_v53.tex}.
\bibitem{Unit}\emph{A single norm determinant for the 25- and 13-variable constructions},
repository manuscript,
\nolinkurl{papers/finite_algebra/unit_norm_20261114_v28.tex}.
\bibitem{Baseline}\emph{An explicit polynomial in 29 variables over the rationals},
repository manuscript, elementary Lucas growth lemma,
\nolinkurl{papers/few_variables/few_variables_20261108_v8.tex}.
\end{thebibliography}
\end{document}
